\section{Host fields, Frobenius classes and the balance identity}\label{sec:hosts}

\subsection{Frames and hosts}

\begin{definition}[Frames]\label{def:frame}
A \emph{frame} $\varphi$ consists of integers $e\ge0$ and $\mathrm{ord}\ge1$, an abelian group $A$ of
order $\mathrm{ord}$, a surjection $h:A\to H_e=(\Z/2^{e+2}\Z)^\times/\{\pm1\}$ and a cyclic subgroup
$A_0\subset A$ on which $h$ is injective. Its index is $o(\varphi)=[A:A_0]$ and its Galois group is
$\mathcal G_\varphi=\{\pm1\}\times\{\pm1\}\times A$. \formal{Formalism.HC.Frame}
\end{definition}

\begin{definition}[Hosts]\label{def:host}
A \emph{host} on a frame $\varphi$ consists of:
\begin{itemize}
\item primes $q\equiv3\pmod8$ and $d\equiv7\pmod{16}$ with $(\frac{-d}{q})=1$;
\item an integer $Q_S$ with $2^{e+2}\mid Q_S$, $\gcd(Q_S,qd)=1$ and $(\frac{-d}{l})=1$ for all odd
$l\mid Q_S$;
\item a Galois number field $E\subset\Q(\zeta_{N_H})$, where $N_H=Q_Sqd$, with an isomorphism
$\Gal(E/\Q)\cong\mathcal G_\varphi$, under which the second factor $\{\pm1\}$ acts on $\sqrt{-d}\in E$;
\item an Artin map $\mathrm{art}:(\Z/N_H)^\times\to\mathcal G_\varphi$, compatible with the cyclotomic
action, with $\mathrm{art}(-1)$ the complex conjugation and compatible with $h$ on the $2$-power part.
\end{itemize}
The host is \emph{cut out by} $\psi\in\mathcal P$ if $N_\psi\mid Q_S$ and, for $a$ prime to $N_H$,
$\mathrm{art}(a)=1$ exactly when $a\equiv\pm1\pmod{2^{e+2}}$, $\psi(a)=1$ and
$(\frac aq)=(\frac ad)=1$. \formal{Formalism.HC.HostRaw} \formal{UpperTheory.Top.KerChi}
\end{definition}

For a host cut out by $\psi$, the characters of $\mathcal G_\varphi$ that are trivial on $A_0$ (in the
first factor) correspond to the $4o$ twists $\xi(\frac{\cdot}{q})^i(\frac{\cdot}{d})^j$ with
$\xi\in\langle\chi_\psi\rangle$ and $i,j\in\{0,1\}$. We call those with $(i,j)=(0,0)$ \emph{kept} and the
others \emph{twisted}. The Frobenius classes in $\mathcal G_\varphi$ of the primes $p\nmid N_H$ refine
their splitting in $\Q(\sqrt{-d})$: $p$ is inert in $\Q(\sqrt{-d})$ exactly when the sign of
$\mathrm{art}(p)$ on $\sqrt{-d}$ is $-1$.

\subsection{The balance identity}

Let $\mathcal E=(\Z/2)^2$ and let $H$ be a finite abelian group (later $H=A$). Put
\[
\Phi=\{((\epsilon_1,\epsilon_2),h):\epsilon_1=0\}\subset\mathcal E\times H,\qquad
c=((1,1),0),\qquad \mathrm{partner}(t,g)=c+g+((0,0),t),
\]
and let $\chi_C$ be the character $((\epsilon_1,\epsilon_2),h)\mapsto(-1)^{\epsilon_1}$.

\begin{definition}[Arithmetic data]\label{def:hosts:arith}
An \emph{arithmetic datum} $D$ on $H$ consists of primes $q_1,d$ and, for each prime $p$, a subgroup
$I_p\le\mathcal E\times H$ and an element $\Fr_p\in\mathcal E\times H$. For a character $\psi$ of
$\mathcal E\times H$ put $\psi^D(n)=\prod_{p\mid n}|I_p|^{-1}\sum_{x\in I_p}\psi(v_p(n)\Fr_p+x)$. Given a
function $\ell$ on characters, put $\mathcal G(\psi,y)=G(y)$ if $\psi=0$, and otherwise
\[
\mathcal G(\psi,y)=\sum_{n\le y}\Lambda(n)\psi^D(n)\Big(\frac yn-1\Big)+y\,\ell(\psi).
\]
A prime $p\notin\{2,q_1,d\}$ is \emph{inert} if $\Fr_p\notin\Phi$. For $N\in\N$, $a>0$, $t\in H$ and a
weight $w:\mathcal E\times H\to\R_{>0}$ define the \emph{inert count}, the \emph{main term} and the
\emph{character average}
\begin{align*}
\mathrm{ASK}(D,N,a,t,w)&=\sum_{\substack{p\le N\\ \text{inert}}}2\log p\sum_{\substack{1\le i\le N\\ i\Fr_p\notin\Phi}}
\Big(\frac a{p^i}-w(\mathrm{partner}(t,i\Fr_p))\Big)^+\\
&\qquad+(\text{the analogous term for }p=d),\\
\mathrm{MAIN}(\ell,a,w)&=a\Big(\log a-1-\gamma+\ell(\chi_C)-\frac1{|\Phi|}\sum_{\rho\in\Phi}\log w(\rho)\Big),\\
\mathrm{De}(D,\ell,a,t,w)&=\frac1{|\Phi|}\sum_{g\notin\Phi}\sum_\psi\overline{\psi(g)}\,w(\mathrm{partner}(t,g))\,
\mathcal G\big(\psi,a/w(\mathrm{partner}(t,g))\big).
\end{align*}
\formal{Formalism.LFunctionChainParameters.characterSumData.ASK}
\end{definition}

For the datum $\mathrm{arith}(H)$ of a host (with $\Fr_p$ the image of $\mathrm{art}(p)$ and $I_p$ the
inertia groups), the inert primes are exactly the primes inert in $\Q(\sqrt{-d})$. In that case
$\mathcal G(\psi,y)=g_{\chi}(y)$, where $\chi$ is the primitive Dirichlet character attached to $\psi$
via $\mathrm{art}$, provided $\ell$ is compatible.

\begin{proposition}[Balance identity]\label{prop:balance}
Assume $D$ satisfies the \emph{host facts}: $I_p=0$ for $p\notin S\cup\{d\}$; for $p\in S$,
$I_p\subset\Phi$ and $\Fr_p\in\Phi$; and $I_d=\{0,((1,0),0)\}$. Here $S$ is a finite set of primes with
$2,q_1\in S\not\ni d$. If $a>0$, $w>0$ on $\Phi$ and $a/w\le N$ on $\Phi$, then for every $\ell$ and
every $t\in H$
\[
\mathrm{ASK}(D,N,a,t,w)-\mathrm{MAIN}(\ell,a,w)=\mathrm{De}(D,\ell,a,t,w).
\]
\formal{Deduction.ch_balance}
\end{proposition}

\begin{proof}
For $\psi\neq0$ and $w>0$,
$\sum_{n\le N}\Lambda(n)\psi^D(n)\max(a/n-w,0)=w\,\mathcal G(\psi,a/w)-a\,\ell(\psi)$. For $\psi=0$ the same
sum equals $w\,\mathcal G(0,a/w)+a\log a-a\log w-(1+\gamma)a$. Insert these into $\mathrm{De}$ and
evaluate the prime-power sums by orthogonality:
\[
\frac1{|\Phi|}\sum_\psi\overline{\psi(g)}\psi^D(p^i)=\begin{cases}2/|I_p|,& g-i\Fr_p\in I_p,\\ 0,&\text{otherwise}\end{cases}
\]
(note $|\mathcal E\times H|=2|\Phi|$). Thus an unramified $p^i$ contributes $2\log p$ at the single
class $g=i\Fr_p$, and only if $i\Fr_p\notin\Phi$, which forces $p$ to be inert. The primes in $S$
contribute nothing, because $i\Fr_p+I_p\subset\Phi$, and $p=d$ gives the second term of
$\mathrm{ASK}$. The remaining terms are summed with
$\sum_{g\notin\Phi}\overline{\psi(g)}=|\Phi|,-|\Phi|,0$ for $\psi=0$, $\psi=\chi_C$ and otherwise, and
with $\sum_{g\notin\Phi}\log w(\mathrm{partner}(t,g))=\sum_{\rho\in\Phi}\log w(\rho)$. They give
$-\mathrm{MAIN}$.
\end{proof}

The identity is an exact explicit formula for inert prime powers in Frobenius classes, with Riesz
weights whose scale $a/w$ depends on the class.

\subsection{Coset distributions}

For $\gamma\in A$ put $\lambda(\gamma)=\frac4{3o+1}$ if $\gamma\in A_0$ and $\frac3{3o+1}$ otherwise.
A finite set $T\subset A$ has \emph{shares within $\delta$} if
$|\#\{t\in T:t-\gamma\in A_0\}-\lambda(\gamma)|T||\le\delta|T|$ for every $\gamma$.

\begin{remark}\label{rem:hosts:fourier}
If $T$ has shares within $\delta$ and $\xi$ is a non-trivial character of the cyclic group $A/A_0$ of
order $o$, then
\[
\sum_{t\in T}\xi(t)=|T|\Big(\frac{4-3}{3o+1}+\frac3{3o+1}\sum_{\gamma\in A/A_0}\xi(\gamma)\Big)+O(o\delta|T|)
=\frac{|T|}{3o+1}+O(o\delta|T|).
\]
Averaging over $T$ therefore produces exactly the weights of \cref{def:landau:Gm}.
\end{remark}

\subsection{The Chebyshev inequality at a host}

\begin{definition}[Host conditions]\label{def:hosts:EAtH}
Let $\psi\in\mathcal P$, $x>1$, $\eps>0$ and $C>0$. A host $H$ cut out by $\psi$, together with a finite
set $T\subset A$, satisfies the \emph{host conditions at} $(x,\eps,C)$ if:
\begin{enumerate}[label=(\roman*)]
\item $o(\varphi)=\mathrm{ord}(\chi_\psi)$;
\item $T$ has shares within $x^{-1+\eps}$;
\item for every twisted character $\theta$, the attached primitive character $\chi$ satisfies
$|g_\chi(x)|\le x^{1/2+\eps}$ and $|\ell(\chi)|\le C\log x$.
\end{enumerate}
\formal{Formalism.HC.EAtH}
\end{definition}

Condition (iii) is a square-root bound at the single point $x$ for the $3o$ twisted characters. It does
not follow from known results for an arbitrary host. It is achieved by the choice of the auxiliary
primes $q,d$ (\cref{thm:supply}).

The arithmetic input at a host is a \emph{gap bound} for the inert counts. Let $j$ run over a set of
$m$ indices, and $u$ over $A_0$. Let $a:T\to\R$ (with $|a(t)/b-x|\le1$) and weights
$w(u,j,\cdot)>0$ be given, with $w\in[b,b+\eta+1]$ on the ``thick'' classes
$\mathrm{partner}(t,\cdot)\in T_2+u$, $T_2\subset A_0$, and $w\le2b'$ elsewhere. The \emph{gap} is
\[
\mathrm{gap}(u)=\sum_{t\in T}\sum_{j}\Big(\mathrm{ASK}\big(\mathrm{arith}(H),N,a(t),t,w(u,j,\cdot)\big)
-\Re\,\mathrm{MAIN}\big(\ell,a(t),w(u,j,\cdot)\big)\Big).
\]

\begin{theorem}[Chebyshev inequality]\label{thm:cheb-windows}
For every $\psi\in\mathcal P$ and all $C,C'\ge0$, $\eps>0$ there is $x_1$ with the following property.
Let $x\ge x_1$, let a host $H$ cut out by $\psi$ and $T\subset A$ satisfy the host conditions at
$(x,\eps,C)$, and suppose that the size conditions below hold. Suppose moreover that
\[
\mathrm{gap}(u)\le x^{C'\eps}\,\mathrm{EB},\qquad
\mathrm{EB}=|T|\,a_E\Big(1+\frac{m^2}k\Big)+m\Big(2nmbx^\eps+\frac{|T_2|kbx^\eps}{n^2x^{-\eps}}\Big)
\]
for all $u\in A_0$, where $n=|A|$ and $a_E\le2xb$. Then
\[
G_\psi(x)\le x^{\frac12+(C'+4)\eps}.
\]
The size conditions are: $|T|\ge x^{-\eps/8}n/2$, $|T_2|\ge x^{-\eps}n$, $x^{1/2-\eps}\le m\le
x^{1/2+\eps/2}$, $xm\le k\le n^2x^{1/2-\eps}$; and the prime sums $\mathcal G(\psi',y)$ are at most
$y\theta$ in absolute value for $y\ge X'/3$, with $\theta$ and $\eta$ small in terms of $x$, $b$ and
$|\Phi|$. \formal{UpperTheory.Comparison₂.core_gm} \formal{Deduction.ch_core}
\end{theorem}

\begin{proof}
Write $\Psi$ for the $4o$ characters of $\mathcal E\times A$ that are trivial on $0\times A_0$. Let $K$
be the kept ones (those attached to the $\psi^k$) and put $c=(3o+1)^{-1}$. Put
\[
\mathrm{GmA}(y)=\sum_{\psi'\in K}\mathrm{wt}(\psi')\Re\mathcal G(\psi',y),\qquad
\mathrm{SigA}(y)=\sum_{\psi'\in\Psi\smallsetminus K}|\mathcal G(\psi',y)|,
\]
with weight $1$ on characters trivial on $0\times A$ and $c$ otherwise. By the compatibility of
$\mathcal G$ with $g_\chi$, $\mathrm{GmA}=G_\psi$. By condition (iii), $\mathrm{SigA}(x)\le4o\,x^{1/2+\eps}$.

(1) \emph{Regularity.} $|\mathcal G(\psi',y)-\mathcal G(\psi',y')|\le|y-y'|(\log\max(y,y')+C_d+|\ell(\psi')|)$
for $\psi'\ne0$, and $|G(y)-G(y')|\le C_d|y-y'|$.

(2) \emph{Equidistribution.} Put $y_t=a(t)/b$, so $|y_t-x|\le1$, and $e_1(t)=((1,0),t)$. Let
$v(\psi')=\mathrm{wt}(\psi')\Re\mathcal G(\psi',x)$ for $\psi'\in K$ and $v(\psi')=-|\mathcal G(\psi',x)|$
otherwise. By (1), \cref{rem:hosts:fourier} and condition (ii),
\[
|T|\big(v(\psi')-A_1x^{1/4+\eps}\big)\le\sum_{t\in T}\Re\big(\psi'(e_1(t))\mathcal G(\psi',y_t)\big).
\]
Summing over $\Psi$,
\[
|T|\,Y\le\sum_t\Re\sum_{\psi'\in\Psi}\psi'(e_1(t))\mathcal G(\psi',y_t),\qquad
Y=\mathrm{GmA}(x)-\mathrm{SigA}(x)-4x^{1/2+\eps}.
\]

(3) \emph{Balance.} Average over $u\in A_0$, using orthogonality on $A_0$ and $T_2\subset A_0$. This
identifies the right-hand side with $\frac{2n}{|A_0|b|T_2|}$ times a sum of the character averages
$\mathrm{De}$ over $t$, $j$ and $u$, up to an error controlled by the size conditions. On the thick
classes the weights $w\in[b,b+\eta+1]$ are compared with $b$ using (1); on the thin classes
$a(t)/w\ge X'/3$, and the prime sums are small by assumption. By \cref{prop:balance},
$\Re\,\mathrm{De}=\mathrm{ASK}-\Re\,\mathrm{MAIN}$, so the gap bound applies and gives
\[
|T|\,m\,b\,\frac{|T_2|}{2n}\,Y\le x^{C'\eps}\mathrm{EB}+\frac{2|T|mb}x.
\]

(4) \emph{Sizes.} Dividing by $|T|mb|T_2|/(2n)$ and using the size conditions, every term is
$O(x^{1/2+3\eps})$. Hence
\[
G_\psi(x)\le x^{1/2+(C'+3)\eps}+\mathrm{SigA}(x)\le x^{1/2+(C'+4)\eps}
\]
for $x$ large.
\end{proof}

The theorem is the only place where the arithmetic input, the gap bound, meets the prime sums $G_\psi$.
In the formal development the gap bound is part of the predicate $\sfname{ChainTied}$
(\cref{def:chaintied}), and the size conditions are part of the requirements on families
(\cref{sec:top}).
